\chapter{Prior-weighted logistic regression analysis}

In this section, we analyze the performance of the Prior-weighted Logistic Regression model. Unlike the standard model, this version explicitly incorporates the target application prior $\pi_T = 0.1$ into the objective function during the training phase. 

\section{Experimental results}

The model was trained on the full project dataset using 13 values of $\lambda$ logarithmically spaced between $10^{-4}$ and $10^{2}$. To obtain the Log-Likelihood Ratios (LLRs), we subtracted the log-odds of the target prior used during training ($\pi_T = 0.1$) instead of the empirical prior.

The numerical results are summarized in Table \ref{tab:weighted_lr_results} and plotted in Figure \ref{fig:lr_weighted_plot}.

\begin{table}[h]
\centering
\caption{Prior-weighted logistic regression results ($\pi_T = 0.1$)}
\label{tab:weighted_lr_results}
\begin{tabular}{@{}ccc@{}}
\toprule
\textbf{Regularization Coefficient ($\lambda$)} & \textbf{minDCF} & \textbf{actDCF} \\ \midrule
$1.0000 \times 10^{-4}$ & 0.3721 & 0.4071 \\
$3.1623 \times 10^{-4}$ & 0.3701 & 0.4010 \\
$1.0000 \times 10^{-3}$ & 0.3699 & 0.4129 \\
$3.1623 \times 10^{-3}$ & 0.3670 & 0.4327 \\
$1.0000 \times 10^{-2}$ & 0.3630 & 0.4487 \\
$3.1623 \times 10^{-2}$ & 0.3640 & 0.5964 \\
$1.0000 \times 10^{-1}$ & 0.3648 & 0.9147 \\
$3.1623 \times 10^{-1}$ & 0.3640 & 1.0000 \\
$1.0000 \times 10^{0}$  & 0.3640 & 1.0000 \\
$3.1623 \times 10^{0}$  & 0.3630 & 1.0000 \\
$1.0000 \times 10^{1}$  & 0.3630 & 1.0000 \\
$3.1623 \times 10^{1}$  & \textbf{0.3620} & 1.0000 \\
$1.0000 \times 10^{2}$  & \textbf{0.3620} & 1.0000 \\
\bottomrule
\end{tabular}
\end{table}

\begin{figure}[h]
    \centering
    \includegraphics[width=0.85\textwidth]{LR_Weighted_DCF_vs_Lambda.png}
    \caption{Actual DCF and Minimum DCF as a function of $\lambda$ for the Prior-weighted Logistic Regression model ($\pi_T = 0.1$).}
    \label{fig:lr_weighted_plot}
\end{figure}

The optimal discriminative capability is reached at high regularization ($\lambda = 3.16 \times 10^{1}$), yielding a minDCF of \textbf{0.3620}. The best calibration is observed at $\lambda = 3.16 \times 10^{-4}$, where actDCF hits \textbf{0.4010}.

\section{Analysis and discussion}

\subsection*{Are there significant differences for this task?}

No, there are no significant differences in terms of performance between the standard and the prior-weighted models for this specific task. Comparing the best minimum DCF obtained by the standard model ($0.3611$) with the best minimum DCF of the prior-weighted model ($0.3620$), the values are nearly identical. Both models exhibit the exact same trend: a stable minDCF across the lambda range and a degrading actDCF that rapidly reaches $1.0$ as regularization increases past $10^{-1}$.

The lack of difference is motivated by the fact that the project dataset is sufficiently large. In high-data regimes, the optimization algorithm has enough information to learn the underlying class distributions and position the separating hyperplane effectively, regardless of whether the minority class errors are explicitly over-weighted in the objective function. The post-hoc threshold adjustment (removing empirical prior log-odds) used in the standard model is mathematically sufficient to recover the optimal Bayes decisions.

\subsection*{Are there advantages using the prior-weighted model for our application?}
\textit{(Remember that the prior-weighted model requires that we know the target prior when we build the model).}

In this specific case, the prior-weighted model does not offer a numerical advantage, but it provides a distinct theoretical one. The primary advantage of the prior-weighted model is its ability to "focus" the training on a specific operating point. By incorporating $\pi_T$ directly into the loss function, we instruct the optimizer to penalize classification errors on the expected minority class more heavily (in our case, the "Genuine" class since $\pi_T=0.1$ implies a 90\% fake distribution). 

Theoretically, the prior-weighted model is most advantageous when the training set is highly imbalanced or completely unrepresentative of the real-world application, allowing the model to adapt its internal weights natively to the deployment scenario.

However, this approach introduces two major practical drawbacks:
\begin{itemize}
    \item \textbf{Requirement of prior knowledge:} We must know the exact target application prior at training time. If the operational requirements change (e.g., if the system administrator later decides to evaluate for a balanced prior of $\pi_T = 0.5$), the prior-weighted model trained specifically for $0.1$ will be poorly calibrated for the new threshold and will yield sub-optimal results.
    \item \textbf{Calibration vulnerability:} As seen in the results, even with the "correct" prior explicitly weighted into the training loop, the actDCF still massively diverges from the minDCF at high $\lambda$ values. This confirms that aggressive $L_2$ regularization still compresses the scores toward the bias term, leading to a total loss of probabilistic meaning that the prior weighting alone cannot salvage.
\end{itemize}